%%%%%%%%%%%%%%%%%%%%%%%%%%%%%%%%%%%%%%%%%%%%%%%%%%%%%%%
%% Engineer & Master Thesis, LaTeX Template          %%
%% Copyleft by Piotr Woźniak & Artur M. Brodzki      %%
%% Faculty of Electronics and Information Technology %%
%% Warsaw University of Technology, Warsaw, 2019     %%
%%%%%%%%%%%%%%%%%%%%%%%%%%%%%%%%%%%%%%%%%%%%%%%%%%%%%%%

\documentclass[
    left=2.5cm,         % Sadly, generic margin parameter
    right=2.5cm,        % doesnt't work, as it is
    top=2.5cm,          % superseded by more specific
    bottom=3cm,         % left...bottom parameters.
    bindingoffset=6mm,  % Optional binding offset.
    nohyphenation=false % You may turn off hyphenation, if don't like.
]{eiti/eiti-thesis}

\langpol % Dla języka angielskiego mamy \langeng
\graphicspath{{img/}}             % Katalog z obrazkami.
\addbibresource{bibliografia.bib} % Plik .bib z bibliografią

\begin{document}

%--------------------------------------
% Strona tytułowa
%--------------------------------------
\MasterThesis % dla pracy inżynierskiej mamy \EngineerThesis
\instytut{XXXXXX}
\kierunek{XXXXXX}
\specjalnosc{XXXXXX}
\title{
    Niepotrzebnie długi i skomplikowany tytuł pracy \\
    trudny do przeczytania, zrozumienia i wymówienia
}
\engtitle{ % Tytuł po angielsku do angielskiego streszczenia
    Unnecessarily long and complicated thesis' title \\
    difficult to read, understand and pronounce
}
\author{\{Imię i Nazwisko\}}
\album{XXXXXX}
\promotor{XXXXXX}
\date{\the\year}
\maketitle

%--------------------------------------
% Streszczenie po polsku
%--------------------------------------
\cleardoublepage % Zaczynamy od nieparzystej strony
\streszczenie \lipsum[1-3]
\slowakluczowe XXX, XXX, XXX

%--------------------------------------
% Streszczenie po angielsku
%--------------------------------------
\newpage
\abstract \kant[1-3]
\keywords XXX, XXX, XXX

%--------------------------------------
% Oświadczenie o autorstwie
%--------------------------------------
\newpage
\makeauthorship

%--------------------------------------
% Spis treści
%--------------------------------------
\cleardoublepage % Zaczynamy od nieparzystej strony
\pagestyle{plain}
\tableofcontents

%--------------------------------------
% Rozdziały
%--------------------------------------
\cleardoublepage % Zaczynamy od nieparzystej strony
\pagestyle{headings}

\newpage % Rozdziały zaczynamy od nowej strony.
\section{Praefatio}
\lipsum[1] \cite{goossens93}
\begin{figure}[!h]
    \label{fig:anzelm}
    \centering \includegraphics[width=0.5\linewidth]{logopw.png}
    \caption{Tradycyjne godło Politechniki Warszawskiej.}
\end{figure}
\lipsum[2-10]
         % Wygodnie jest trzymać każdy rozdział w osobnym pliku.
\newpage % Rozdziały zaczynamy od nowej strony.
\section{De Finibus Bonorum et Malorum}
\lipsum[1] Lorem ipsum dolor sit amet\footnote{Lorem ipsum dolor sit amet, consectetur adipiscing elit, sed do eiusmod tempor incididunt ut labore et dolore magna aliqua. Ut enim ad minim veniam, quis nostrud exercitation ullamco laboris nisi ut aliquip ex ea commodo consequat.}. 
\begin{align*}
    E & = mc^2 \\ 
    y & = ax^2 + bx + c
\end{align*}

\lipsum[3]

\begin{align}
\begin{bmatrix}
    1 & 0 & 0 \\ 
    0 & 2 & 0 \\ 
    0 & 0 & 3
\end{bmatrix} \cdot 
\begin{bmatrix}
    4 \\ 5 \\ 6
\end{bmatrix} = 
\begin{bmatrix}
    4 \\ 10 \\ 18
\end{bmatrix}
\end{align}

\lipsum[4] Lorem ipsum dolor sit amet, consectetur adipiscing elit, sed do eiusmod tempor incididunt ut labore et dolore magna aliqua \cite{szczypiorski2015}, \cite{duqu2011}, \cite{shs2015}, \cite{wozniak2018}, \cite{dcp19}. 

\subsection{Critique of Pure Reason}
\kant[1]

\begin{table}[!h] \label{tab:tabela1} \centering
\caption{Przykładowa tabela.}
\begin{tabular} {| c | c | r |} \hline
    Kolumna 1 & Kolumna 2 & Liczba \\ \hline\hline
    cell1 & cell2 & 60 \\ \hline
    cell4 & cell5 & 43 \\ \hline
    cell7 & cell8 & 20,45 \\ \hline
    \multicolumn{2}{|r|}{Suma:} & 123,45 \\ \hline
\end{tabular}
\end{table}

\kant[2]

\begin{longtable}{| c | m{0.58\linewidth} | r | m{0.1\linewidth} |} 
    \caption{Tabela wielostronicowa.} \\ 
    \hline
    Lp & \multicolumn{1}{c|}{Treść} & \multicolumn{1}{c|}{Kwota} & \multicolumn{1}{m{0.1\linewidth}|}{Wariant opłaty} \\ \hline\hline \endfirsthead
    
    \endfoot
    \hline \endlastfoot
    
    1 & Lorem ipsum dolor sit amet, consectetur adipiscing elit, sed do eiusmod tempor incididunt ut labore et dolore magna aliqua. & 111 111,11 zł & \multicolumn{1}{c|}{WAR1} \\ \hline
    2 & Lorem ipsum dolor sit amet, consectetur adipiscing elit, sed do eiusmod tempor incididunt ut labore et dolore magna aliqua. & 22 222,22 zł & \multicolumn{1}{c|}{WAR1} \\ \hline
    3 & Lorem ipsum dolor sit amet, consectetur adipiscing elit, sed do eiusmod tempor incididunt ut labore et dolore magna aliqua. & 33 333,33 zł & \multicolumn{1}{c|}{WAR1} \\ \hline
    4 & Lorem ipsum dolor sit amet, consectetur adipiscing elit, sed do eiusmod tempor incididunt ut labore et dolore magna aliqua. & 444 444,44 zł & \multicolumn{1}{c|}{WAR1} \\ \hline
    5 & Lorem ipsum dolor sit amet, consectetur adipiscing elit, sed do eiusmod tempor incididunt ut labore et dolore magna aliqua. & 55 555,55 zł & \multicolumn{1}{c|}{WAR1} \\ \hline
    6 & Lorem ipsum dolor sit amet, consectetur adipiscing elit, sed do eiusmod tempor incididunt ut labore et dolore magna aliqua. & 66 666,66 zł & \multicolumn{1}{c|}{WAR1} \\ \hline
    7 & Lorem ipsum dolor sit amet, consectetur adipiscing elit, sed do eiusmod tempor incididunt ut labore et dolore magna aliqua. & 777 777,77 zł & \multicolumn{1}{c|}{WAR1} \\ \hline
    8 & Lorem ipsum dolor sit amet, consectetur adipiscing elit, sed do eiusmod tempor incididunt ut labore et dolore magna aliqua. & 8 888,88 zł & \multicolumn{1}{c|}{WAR1} \\ \hline
    9 & Lorem ipsum dolor sit amet, consectetur adipiscing elit, sed do eiusmod tempor incididunt ut labore et dolore magna aliqua. & 999 999,99 zł & \multicolumn{1}{c|}{WAR1} \\ \hline
    10 & Lorem ipsum dolor sit amet, consectetur adipiscing elit, sed do eiusmod tempor incididunt ut labore et dolore magna aliqua. & 111 111,11 zł & \multicolumn{1}{c|}{WAR2} \\ \hline
    11 & Lorem ipsum dolor sit amet, consectetur adipiscing elit, sed do eiusmod tempor incididunt ut labore et dolore magna aliqua. & 22 222,22 zł & \multicolumn{1}{c|}{WAR2} \\ \hline
    12 & Lorem ipsum dolor sit amet, consectetur adipiscing elit, sed do eiusmod tempor incididunt ut labore et dolore magna aliqua. & 33 333,33 zł & \multicolumn{1}{c|}{WAR2} \\ \hline
    13 & Lorem ipsum dolor sit amet, consectetur adipiscing elit, sed do eiusmod tempor incididunt ut labore et dolore magna aliqua. & 444 444,44 zł & \multicolumn{1}{c|}{WAR2} \\ \hline
    14 & Lorem ipsum dolor sit amet, consectetur adipiscing elit, sed do eiusmod tempor incididunt ut labore et dolore magna aliqua. & 55 555,55 zł & \multicolumn{1}{c|}{WAR2} \\ \hline
    15 & Lorem ipsum dolor sit amet, consectetur adipiscing elit, sed do eiusmod tempor incididunt ut labore et dolore magna aliqua. & 66 666,66 zł & \multicolumn{1}{c|}{WAR2} \\ \hline
    & \multicolumn{1}{r|}{\textbf{Suma:}} & \textbf{7 777 777,77 zł} & 
    \label{table:koszty}
\end{longtable}
\kant[4]

\subsection{Categorical Imperative}
As any dedicated reader can clearly see, the Ideal of practical reason is a representation of, as far as I know, the things in themselves; as I have shown elsewhere, the phenomena should only be used as a canon for our understanding:
% Parametr label ustawia symbol, a leftmargin - wielkość wcięcia.
% Domyślny układ to [---] bez wcięcia, bo tak pan Marcin Woliński powiedział;
% ale ja nie polecam. // AB
\begin{itemize}
    \item Item 1:
    \begin{itemize}[label=---]
        \item item 1.1;
        \item item 1.2;
        \item item 1.3;
    \end{itemize}
    \item Item 2;
    \item Item 3;
    \item Item 4.
\end{itemize}
\kant[2]
\begin{enumerate}
    \item Item 1:
    \begin{enumerate}
        \item item 1.1;
        \item item 1.2:
        \begin{enumerate}
            \item item 1.2.1;
            \item item 1.2.2;
        \end{enumerate}
        \item item 1.3;
    \end{enumerate}
    \item Item 2;
    \item Item 3;
    \item Item 4.
\end{enumerate}

\kant[9]

\subsection{G\"odel's ontological proof}
\kant[9] Lorem ipsum dolor sit amet, consectetur adipiscing elit, sed do eiusmod tempor incididunt ut labore et dolore magna aliqua \cite{benzmuller2014}, \cite{goedel95}, \cite{wang97}, \cite{koons2005}. 
\begin{assumption} \label{ass:1}
    $ [\![ \ \phi \ ]\!] \Longrightarrow [\![ \ P(\phi); \neg P(\phi) \ ]\!]$
\end{assumption}
\begin{axiom}[Dualność] \label{axiom:1}
    $\neg P(\phi) \Leftrightarrow P(\neg \phi)$, równoważnie $P(\phi) \Leftrightarrow \neg P(\neg \phi)$
\end{axiom}
\begin{axiom}[Całkowitość] \label{axiom:2}
    $ \left( P(\phi) \wedge \forall x: \phi(x) \Rightarrow \psi(x) \right) \Rightarrow P(\psi) $
\end{axiom}
\begin{axiom}[Absolutność] \label{axiom:3}
    $ P(\phi) \Rightarrow \Box P(\phi) $
\end{axiom}
\begin{definition} \label{def:1}
    $ G(x) \Leftrightarrow \forall \phi: \left( P(\phi) \Rightarrow \phi(x) \right) $
\end{definition}
\begin{definition} \label{def:2}
    $ \phi \ ess \ x \Leftrightarrow \phi(x) \wedge \forall \psi \left( \psi(x) \Rightarrow \Box \forall y \left( \phi(y) \Rightarrow \psi(y) \right) \right)  $
\end{definition}
\begin{axiom} \label{axiom:4}
    P(G)
\end{axiom}
\begin{lemma} \label{lemma:1}
    $ P(\phi) \Rightarrow \Diamond \exists x : \phi(x) $
\end{lemma}
\begin{proof}
    Dowód pomijamy, bo jest trywialny :)
\end{proof}
\begin{lemma} \label{lemma:2}
    $ \Diamond \exists x : G(x) $
\end{lemma}
\begin{proof}
    Natychmiastowy wniosek z aksjomatu \ref{axiom:4} i lematu \ref{lemma:1}.
\end{proof}
\begin{lemma} \label{lemma:3}
    $ G(x) \Rightarrow G \ ess \ x $
\end{lemma}
\begin{proof}
    Poprzez podstawienie do definicji \ref{def:2}. 
\end{proof}
\begin{definition} \label{def:3}
    $ E(x) \Leftrightarrow \forall \phi \left( \phi \ ess \ x \Rightarrow \Box\ \exists x: \phi(x) \right) $
\end{definition}
\begin{axiom} \label{axiom:5}
    P(E)
\end{axiom}
\begin{theorem}
    $ \Box\ \exists x : G(x) $
\end{theorem}
\begin{proof}
    Na podstawie definicji \ref{def:1}, lematu \ref{lemma:3} i aksjomatu \ref{axiom:5}. 
\end{proof}    % Umożliwia to również łatwą migrację do nowej wersji szablonu:
\newpage % Rozdziały zaczynamy od nowej strony.
\section{Code listings}
\lipsum[10]
\begin{lstlisting}[language=HTML]
<html>
  <head>
    <title>Hello world!</title>
  </head>
  <body>
    Hello
  </body>
</html>
\end{lstlisting}
\lipsum[11]
\begin{lstlisting}[language=C]
#include <stdio.h>
int main() {
  // printf() displays the string inside quotation
  printf("Hello world!");
  return 0;
}
\end{lstlisting}
\lipsum[12]
 % wystarczy podmienić swoje pliki main.tex i eiti-thesis.cls
                            % na nowe wersje, a cały tekst pracy pozostaje nienaruszony.
\newpage % Rozdziały zaczynamy od nowej strony.
\section{Zbiory danych i ich wstępna obróbka} \label{dane_wejsciowe}
\subsection{Zbiory danych} \label{dane_wejsciowe}

%todo:
% \colorbox{yellow}{todo:}\\

% //todo: opisać te kategrię ironii z uwzględniem kategorii z wczesniejszego rodziału
% // to z biegunowości i bez to są dwa typy ironii verbalnej -> tylko jedno z nich jest trudniejsze z punktu widzenia zagadanienai kalsyfikacji

W ramach pracy zostały wykorzystane dwa zbiory danych. Pierwszy zbiór danych (w pracy oznaczony jako zbiór A) został udostępniony w ramach inicjatywy SemEval. Inicjatywa ta ma na celu rozwój szeroko pojętej analizy i przetwarzania języka naturalnego. Zbiór ten składa się 4618 próbek oznaczonych tagiem klasyfikującym rekord jako ironiczny lub nie. Dane pochodzą z Twittera i reprezentują trzy różne typy ironii:
\begin{itemize}
    \item Słowną werbalną stworzoną poprzez wykorzystanie przeciwnej biegunowości słów (polarity contrast):
          \begin{itemize}
              \item Przykład: I love waking up with migraines \#not :'(
          \end{itemize}

    \item Słowną werbalną stworzoną bez wykorzystania przeciwnej biegunowości:
          \begin{itemize}
              \item Przykład: Human brains disappear every day. Some of them have never even appeared. \#brain \#humanbrain \#Sarcasm
          \end{itemize}

    \item Ironię sytuacyjną:
          \begin{itemize}
              \item Przykład: Most of us didn't focus in the \#ADHD lecture. \#irony
          \end{itemize}
\end{itemize}



Drugi zbiór (w pracy oznaczony jako zbiór B) także pochodzi z Twittera i został zebrany w ramach publikacji zajmującej się analizą różnic między sarkazmem, a ironią  \cite{Ling2016}  , badacze zebrane dane podzielili na 4 zbiory, każdy o zawierający 30 tysięcy próbek,  reprezentujące poniższe kategorie:
\begin{itemize}
    \item Ironiczny
    \item Sarkastyczny
    \item Zmieszany zbiór ironiczny i sarkastyczny
    \item Niezawierający ani ironii, ani sarkazmu

\end{itemize}

Każdy rekord w ramach zbioru danych posiadał następujące informacje:
\begin{itemize}
    \item datę opublikowania
    \item nazwę użytkownika
    \item unikalny ID Tweetu
\end{itemize}


%todo
% \colorbox{yellow}{todo:}\\

% //todo: lista paramterów

Ze względu na to, że wiadomość zawarta w ramach Tweetu nie była podana wprost, tylko sprecyzowany był jej unikalny ID, konieczne było dokonanie dodatkowego mapowania między rekordami, a postami wykorzystując udostępnione przez Tweeter API. Podczas procesu mapowania nie udało się uzyskać wszystkich treści postów, część z nich była nie osiągalna ze względu na różne czynniki losowe takie jak usuniecie konta użytkownika, zablokowanie Tweetu ze względu na naruszenie regulaminu portalu oraz usunięcie Tweetu przez użytkownika. Dlatego uzyskany zbiór postów był mniejszy niż bazowe 30 tysięcy.  
Na potrzeby pracy zostały wykorzystane 2 zbiory z wcześniej wymienionych:
\begin{itemize}
    \item Zmieszany zbiór ironiczny i sarkastyczny -> zbiór posiadał 22402 próbek
    \item Niezawierający ani ironii, ani sarkazmu -> zbiór posiadł 19018 próbek
\end{itemize}


 
\newpage % Rozdziały zaczynamy od nowej strony.
\section{Preprocessing}

Aby zapewnić jak najlepszą jakość klasyfikacji konieczne jest obrobienie danych do postaci pozbawionej zbędnych szumów. Na potrzeby tej pracy zostały podjęte następujące kroki obróbki danych:

%MOżna napisać czyimi pracami się sugerowałem pisząc preprocessing

\begin{enumerate}
    \item Zastąpienie linków tagiem “<url>”.
    \item Usunięcie znaków końca linii.
    \item Zastąpienie wystąpień nazw użytkowników tagiem  “<username>”.
    \item Zamiana emotikon z formatu “:XX\_YY:” na format “<emote> XX YY” .
    \item Rozbicie hasztagów na mniejsze słowa, oznaczenie go tagiem “<hashtag>”.
    \item Zastąpienie oznaczeń odnoszących się do czasu tagiem “<time>”.
    \item Zastąpienie liczb tagiem “<number>”.
    \item Usunięcie znaków  '' | ''  oraz wielokropków.
    \item Konwersja do małych liter.
    \item Rozwinięcie form skróconych do ich pełnych form.
\end{enumerate}

\begin{table}[!h] \label{tab:tabela2} \centering
    \caption{Tabela przedstawiająca efekt preprocessingu.}
    \begin{tabular} {| c | c | c |} \hline
        Numer operacji preprocessingu & Przed konwersją                                          & Po konwersji                   \\ \hline\hline
        1                             & www.google.pl                                            & <url>                          \\ \hline
        2                             & \textbackslash n \quad \textbackslash r \textbackslash n &                                \\ \hline
        3                             & @mike                                                    & <username>                     \\ \hline
        4                             & :face\_screaming\_in\_fear:                              & <emote> face screaming in fear \\ \hline
        5                             & \#trueFriend                                             & <hashtag> true Friend          \\ \hline
        6                             & 14:58                                                    & <time>                         \\ \hline
        7                             & 1234                                                     & <number>                       \\ \hline
        8                             & | ...                                                    & .                              \\ \hline
        9                             & Kazimierz Wielki                                         & kazimierz wielki               \\ \hline
        10                            & I'll \quad | \quad He's                                  & I will \quad | \quad He has    \\ \hline
    \end{tabular}
\end{table}

%todo:
% być może warto opisać w jaki sposób robię te operacje usuwania ?? -> że regexxy
% być może warto opisać o tych formach skróconych, że to model dodakowy wychwytuje
% można teżnapisać czemu nie wykorzystuję stremingu i lematyzacji
%


 
\newpage % Rozdziały zaczynamy od nowej strony.
\section{Eksperymenty}

\colorbox{yellow}{todo: opisać to jakoś sensownie, punkt po punkcie co i dlaczego robię, bo teraz jest tak mega chatycznie, że nie da sie nic wywniskować}\\

//dlaczego takie sieci został ywybran, dlaczeg otakimi paramterami steruję -> na co powinny one wpływać

%W ramach pracy została dokonana analiza wpływu oznaczeń morfosyntaktycznych na jakość klasyfikacji tekstów ironicznych. Na potrzeby języka angielskiego taka analiza sprowadza się głownie do analizy wpływu typu części mowy do jakiego należą poszczególne słowa i taka analiza zostanie podjęta w kolejnych krokach.

% \colorbox{yellow}{highlight}\\

Na potrzeby analizy wykorzystano wiele różnych modeli sieci neuronowych, bazujących głownie na trzech typach warstw:

\begin{itemize}
    \item gęstej;
    \item konwolucyjnej;
    \item LSTM;
\end{itemize}

W ramach różnych architektur testowany był wpływ kilku hiperparametrów na jakość klasyfikacji. Oto lista niektórych z nich:
\begin{itemize}
    \item dla warstwy gęstej:
          \begin{itemize}
              \item  1/2/3 warstwy
              \item Z dropoutem / bez dropoutu
              \item Wielkość wartwy ukrytej 30-200
          \end{itemize}

    \item dla warstwy konwolucyjnej;
          \begin{itemize}
              \item  1/2/3 warstwy
              \item Wielkość okna 3-15
              \item Z dropoutem / bez dropoutu
              \item Wielkość wartwy ukrytej 30-200
          \end{itemize}
    \item dla warstwy LSTM;
          \begin{itemize}
              \item Wykorzystanie warstwy pojedynczej (LSTM) oraz podwójnej (Bi-LSTM)
              \item  1/2/3 warstwy
              \item Z dropoutem / bez dropoutu
              \item Wielkość wartwy ukrytej 30-200
          \end{itemize}
\end{itemize}

% \subsection{POSTAGS}
% %todo:
% //todo: zmergować z dugim rozdiałem o postagach

% Aby możliwe było wykorzystanie przez sieć neuronową informacji na temat typu części mowy danego słowa konieczne są następujące kroki:

% \begin{enumerate}
%     \item Rozpoznanie części mowy w kontekście zdania i przypisanie do niego odpowiedniego tagu,
%     \item Konwersja wykrytego tagu do liczbowej reprezentacji wektorowej,
%     \item Dołączenie utworzonego wektora do przetrenowanego wektora embedingów.
% \end{enumerate}

% \hfill \break % to dodaje pustą linie po czymść dziwnym -lista, tabela
% Do rozpoznania części mowy zostało wykorzystane narzędzie z pakietu NLTK, które pozwala na oznaczenie każdego słowa w zdaniu jednym z 46 tagów, każdy odpowiadający innej części mowy. Następnie wykryty tag został przetworzony do postaci “one-hot encoding” i dodany do istniejącego embedingu słowa, zwiększając tym samy jego rozmiar o długość 46.


\subsection{Wyniki analiza}

\subsubsection{Eksperyment 1}
% Aby w pełni zweryfikować wpływ części mowy (POS tagów) na jakość klasyfikacji zostały wykorzystane trzy różne embeddingi oraz dwa różne zbiory danych wejściowych. Embeddingi różnią się one między innymi długością oraz metodą ich uzyskiwania. A zbiory danych różnią się licznością i stopniem zaszumienia. Każdy ze zbiorów danych został podzielony na trzy zbiory treningowy, walidacyjny oraz testowy w proporcjach 60/20/20 procent.

Aby w pełni zweryfikować wpływ części mowy (POS tagów) na jakość klasyfikacji zostały wykorzystane trzy różne embeddingi oraz dwa różne zbiory danych wejściowych. Wykorzsytane embeddingi to: Glove, FastText oraz ELMo, a ich dokładniejszy opis znajduje się w rozdziale \ref{embeddings}. Natomiast dwa zbiory danych wejściowych (oznaczonych jako A oraz B) są opisane w rozdziale \ref{dane_wejsciowe}. Na potrzeby budowania modeli każdy ze zbiorów danych został podzielony na trzy podzbiory: treningowy, walidacyjny oraz testowy w proporcjach 60/20/20 procent.

% Poniżej przedstawiono tabele z wynikami obliczeń.
Wyniki uzyskane przez modele na danych testowych zbioru A zostały umiszczone w rozdziale \ref{wyniki_eksperymentow} w tabelach \ref{tab:wyniki_glove_A}, \ref{tab:wyniki_fasttext_A}, \ref{tab:wyniki_elmo_A} , a dla zbioru B w tabelach \ref{tab:wyniki_glove_B}, \ref{tab:wyniki_fasttext_B}, \ref{tab:wyniki_elmo_B}.


W oparciu o uzyskane wyniki można stwierdzić, że wraz z dołączeniem informacji o częściach mowy nie obserwuje się widocznej poprawy jakości klasyfikacji, czasem występuje nawet jej pogorszenie. Jednym z powodów takiego zachowania może być fakt, że rodzaj części mowy nie jest kluczowy do rozpoznania ironicznego charakteru tekstu. Możne to być też spowodowane tym, że sieci wychwytują niezbędne dla nich informacje na temat roli konkretnego słowa w zdaniu w oparciu o uzyskiwany kontekst i dodatkowa informacja w ramach embeddingu jest już niepotrzebna.

Ponadto można zaważyć, że wraz ze wzrostem rozmiaru embedingów poprawia się jakość klasyfikacji, co jest zgodne z oczekiwaniami, jako że dłuższy embedding pozwala na zachowanie większej liczby informacji na temat słowa.


Analizując uśrednione pomiary jakości klasyfikacji dla sieci opartych o różne typy warstw, można zauważyć, że najgorzej sprawują się sieci oparte o warstwy gęste. Natomiast sieci wykorzystujące warstwy konwolucyjne oraz LSTM posiadają podobną skuteczność klasyfikacji. Aby wykluczyć wpływ zaszumienia danych oraz przeuczenia sieci wykonano jeszcze dwa eksperymenty mające porównać skuteczność klasyfikacji sieci opartych na warstwach konwolucyjnych oraz LSTM. Eksperymenty te są przeprowadzone na danych wejściowych pozbawionych informacji o częściach mowy, gdyż z dotychczasowych pomiarów wynika, iż mają one pomijalny wpływ na jakość klasyfikacji.

%\colorbox{yellow}{todo: todo opisanie, że zbior duży to zbiór B, a mały zbiór to zbiór A}\\
%todo opisanie, że zbior duży to zbiór B, a mały zbiór to zbiór A

\subsubsection{Eksperyment 2 - Uczenie modelu na zbiorze B i testowanie jakoś klasyfikacji na zbiorze A }

W ramach tego eksperymentu zbiór danych oznaczonych jako \textbf{B} został podzielony na dwa zbiory: treningowy oraz w walidacyjny w proporcjach 80/20 procent. Natomiast zbiór oznaczony jako A został w całości wykorzystany jako zbiór testowy. Wyniki zawierają tabele poniżej.

%todo:
\colorbox{yellow}{todo:}\\
% //tabela dla różnych embedingów pokazująca, że sieci oparte o warstwy konwolucyjne mają najlepszą zdolnośc do generalizacji
//tablea gdzie jest jakość na poziomie 50-55 procent

%todo:
\colorbox{yellow}{todo:}\\
% //Wymyślić wnioski dla teakiej sytuacji
W oparciu o uzyskane miary jakość klasyfikacji można stwierdzić, że modele wyuczone na zbiorze B źle radzą sobie z klasyfikacją rekordów ze zbioru A. Dokładność na poziomie pięćdziesięciu paru procent wskazuje, że pomimo tego, że oba zbiory pochodzą z tego samego źródła to różnią się one na tyle bardzo swoją strukturą wewnętrzną, że żadna architektura nie może sobie poradzić z prawidłową klasyfikacją. Co potwierdza powszechnie znaną prawdę dotyczącą sieci neuronowych, że takie modele do procesu trenowania potrzebują bardzo dużej ilości danych o możliwie jak najbardziej zróżnicowanej strukturze. Wystarczy, że zbiór danych treningowych jest tylko pewnym wycinkiem dziedziny problemu, który sieć ma rozwiązać, by nauczyła się ona klasyfikacji złych cech.


\subsubsection{Eksperyment 3 - Uczenie i testowanie modelu na fragmencie zbioru dużego B tak by był równoliczny ze zbiorem małym A}

W ramach tego eksperymentu ze zbioru danych oznaczonych jako \textbf{B} zostało wybranych w losowy sposób 2000 rekordów ironicznych i tyle samo nieironicznych, tak by licznością były porównywalne ze zbiorem oznaczonym jako \textbf{A}. Tak uzyskany zbiór został podzielony na trzy podzbiory: treningowy, walidacyjny oraz testowy w proporcjach 60/20/20 procent. Wyniki zawierają tabele poniżej.

%todo:
\colorbox{yellow}{todo:}\\
//tabela dla roznych embedingow pokazujaca, że wartwa konwolucyjna jest o 5 procent gorsze od LSTM


W oparciu o uzyskane wyniki można zauważyć:

\begin{itemize}
    \item Wyższą średnią jakoś klasyfikacji w porównaniu do zbioru oznaczonego jako \textbf{A}
    \item Lepszą jakość klasyfikacji dla sieci opartych o warstwy typu LSTM
\end{itemize}


Wyższa jakoś klasyfikacji może wskazywać na to, że zbiór oznaczony jako \textbf{A} posiada dużo bardziej zróżnicowane dane, przy czym zbiór posiada zbyt małą liczbą próbek by sieć, w procesie nauki, odpowiednio je wychwyciła. Problem ten może być najprawdopodobniej wyeliminowany poprzez dostarczenie większej liczby próbek danych o zróżnicowanym charakterze.

Lepsza jakość klasyfikacji przez sieci oparte o warstwy LSTM jest zgodny z wnioskami płynącymi z innych publikacji badających podobne zagadnienia. Trudności z obserwacji tej zależności w poprzednich eksperymentach mogła wynikać między innymi z:


\begin{itemize}

    \item Dla zbioru A:
          \begin{itemize}
              \item Małej liczby próbek
              \item Dużej różnorodności prezentowanej ironii w próbkach
          \end{itemize}

    \item Dla zbioru B:
          \begin{itemize}
              \item Małej różnorodności prezentowanej ironii w próbkach
              \item Na tyle duża liczba próbek, by sieci o teoretycznie gorszych zdolnościach do rozwiązywania danego zadania klasyfikacji mogły wychwycić cechy tekstu pozwalające, przy odpowiednio długim procesie uczenia, na poprawną klasyfikację.
          \end{itemize}

\end{itemize}






 

\newpage % Rozdziały zaczynamy od nowej strony
\section{Summatio}          % Można też pisać rozdziały w jednym pliku.
\lipsum[5-10]

%--------------------------------------------
% Literatura
%--------------------------------------------
\cleardoublepage % Zaczynamy od nieparzystej strony
\printbibliography

%--------------------------------------------
% Spisy (opcjonalne)
%--------------------------------------------
\newpage
\pagestyle{plain}

% Wykaz symboli i skrótów.
% Pamiętaj, żeby posortować symbole alfabetycznie
% we własnym zakresie. Ponieważ mało kto używa takiego wykazu,
% uznałem, że robienie automatycznie sortowanej listy
% na poziomie LaTeXa to za duży overkill.
% Makro \acronymlist generuje właściwy tytuł sekcji,
% w zależności od języka.
% Makro \acronym dodaje skrót/symbol do listy,
% zapewniając podstawowe formatowanie.
% //AB
\vspace{0.8cm}
\acronymlist
\acronym{EiTI}{Wydział Elektroniki i Technik Informacyjnych}
\acronym{PW}{Politechnika Warszawska}
\acronym{WEIRD}{ang. \emph{Western, Educated, Industrialized, Rich and Democratic}}

\listoffigures              % Spis obrazków.
\vspace{1cm}                % vertical space
\listoftables               % Spis tabel.
\vspace{1cm}                % vertical space
\listofappendices           % Spis załączników

% Załączniki
\newpage
\appendix{Nazwa załącznika 1}
\lipsum[1-8]

\newpage
\appendix{Nazwa załącznika 2}
\lipsum[1-4]

% Używając powyższych spisów jako szablonu,
% możesz tu dodać swój własny wykaz bądź listę,
% np. spis algorytmów.

\end{document} % Dobranoc.
